\documentclass[10pt]{article}

\usepackage[utf8]{inputenc}

\usepackage[T1]{fontenc}

\usepackage{amsmath}

\usepackage{amsfonts}

\usepackage{amssymb}

\usepackage[version=4]{mhchem}

\usepackage{stmaryrd}

\usepackage{mathrsfs}



\begin{document}

ствия операторов $a_{\lambda}^{+}$и $\tilde{a}_{\lambda}^{+}$на вакуумное состояние $|0\rangle$. Так, $a_{\lambda}^{+}(p)|0\rangle$ и $\tilde{a}_{\lambda}^{+}(p)|0\rangle$ представляют собой векторы одночастичного состояния (соответственно) частицы и античастицы с импульсом $p$ и спиральностью $\lambda$, а вектор



\[

a_{\lambda_{1}}^{+}\left(p_{1}\right) a_{\lambda_{2}}^{+}\left(p_{2}\right) \ldots a_{\lambda_{n}}^{+}\left(p_{n}\right)|0\rangle

\]



описывает состояние $n$ частиц с импульсами и спиральностями $p_{1}, \lambda_{1} ; p_{2}, \lambda_{2} ; \ldots ; p_{n}, \lambda_{n}$. В соответствии с Ферми-Дирака статистикой вектор (8) антисимметричен относительно перестановки любой пары переменных $p_{i}, \lambda_{i}, i=1,2, \ldots, n$.



Лит.: [1] Боголюбов Н. Н., Ш и рков Д. В., Введение в теорию квантованных полей, 4 изд., М., 1984; [2] Бьёркен Д. Д., Д релл С. Д., Релятивистская квантовая теория, пер. с англ., т. 2, М., $1978 . \quad$ С. М. Биленький. ДИРАКА ПРЕДСТАВЛЕНИЕ - см. Взаимодействия представление.



ДИРАКА СТРУНА - см. Монополь.



ДИРАКА УРАВНЕНИЕ - квантовое (волновое) уравнение для релятивистской частицы со спином 1/2 (электрона, мюона, кварка и других частиц). Получено (для электрона) в 1928 П. Дираком (P. Dirac) из следующих требований:



\begin{enumerate}

  \item уравнение для волновой функции частицы $\psi(\boldsymbol{x}, t)(\boldsymbol{x}-$ пространственные координаты, $t$ - время) должно быть линейным для того, чтобы выполнялся принцип суперпозиции состояний;



  \item в уравнение должна входить первая производная функции $\psi(x, t)$ по времени с тем, чтобы задание $\psi$ в начальный момент определяло волновую функцию в любой последуюший момент времени;



  \item уравнение должно быть инвариантным относительно Лорениа преобразований, то есть иметь один и тот же вид во всех инерциальных системах отсчета;



  \item величина $\psi^{+}(x, t) \psi(x, t)$ (где ${ }^{+}$означает эрмитово сопряжение) должна иметь физич. смысл плотности вероятности нахождения частицы в точке $\boldsymbol{x}$ в момент времени $t$;



  \item уравнение для свободной частицы (массы $m$ ) должно быть построено так, чтобы состояние с импульсом $\boldsymbol{p}$ и энергией \& было его решением только в случае, если выполняется релятивистское соотношение $\mathscr{~}^{2}=\boldsymbol{p}^{2}+m^{2}$ (используется система единиц $\hbar=c=1$ ).



\end{enumerate}



Всем этим требованиям удовлетворяет система уравнений для функции $\psi(x)$, к-рая имеет четыре компоненты и записывается в виде столбца:



\[

\psi(x)=\left(\begin{array}{l}

\psi_{1}(x) \\

\psi_{2}(x) \\

\psi_{3}(x) \\

\psi_{4}(x)

\end{array}\right)

\]



( $x$ - точка пространства-времени). При преобразованиях Лоренца и пространственных поворотах они преобразуются как компоненты 4-компонентного спинора (биспинора).



Ковариантный вид Д. у. зависит от выбора метрики пространства-времени. Если метрика выбрана так, что $x^{2}=g_{\mu v} x^{\mu} x^{\nu}=\left(x^{0}\right)^{2}-x^{2}$, где $g_{\mu \nu}-$ метрич. тензор $\left(x^{0}=t\right)$, то уравнение имеет вид



\[

i \gamma^{\mu} \frac{\partial \psi(x)}{\partial x^{\mu}}-m \psi(x)=0

\]



где $\gamma^{\mu}$ - Дирака матрицы, $\mu=0,1,2,3$ (по повторяющемуся индексу предполагается суммирование). Сопряженный биспинор $\bar{\psi}(x)=\psi^{+}(x) \gamma^{0}$ удовлетворяет уравнению



\[

i \frac{\partial \bar{\psi}(x)}{\partial x^{\mu}} \gamma^{\mu}+\bar{\psi}(x) m=0 .

\]



Из (1) и (2) для 4-мерного вектора тока $j^{\mu}=\bar{\psi} \gamma^{\mu} \psi$ вытекает уравнение непрерывности:



\[

\frac{\partial j^{\mu}}{\partial x^{\mu}}=0

\]



Временна́я компонента вектора тока равна плотности вероятности нахождения частицы в точке $\boldsymbol{x}$ в момент времени $x^{0}$, а его пространственные компоненты являются компонентами 3-мерного вектора потока вероятности.



При данном импульсе $\boldsymbol{p}$ Д. у. имеет четыре линейно независимых решения: два решения с положительной энергией $\mathscr{E}=p_{0}\left(p_{0}=\sqrt{\boldsymbol{p}^{2}+m^{2}}\right)$ и два решения с отрицательной энергией $\mathscr{E}=-p_{0}$. Они могут быть записаны (соответственно) в следующем ковариантном виде:



\[

\Psi_{ \pm p}(x)=\frac{1}{(2 \pi)^{3 / 2}} u( \pm p) e^{\mp i p x},

\]



где спиноры $u(p), u(-p)$ удовлетворяют уравнениям



\[

(\hat{p} \mp m) u( \pm p)=0

\]



$\left(\hat{p}=\gamma^{\mu} p_{\mu}=\gamma^{0} p^{0}-\gamma^{\alpha} p^{\alpha}, \alpha=1,2,3\right)$.



Для сопряженных спиноров имеем:



\[

\bar{u}( \pm p)(\hat{p} \mp m)=0 .

\]



Для каждой из пар спиноров в качестве независимых могут быть выбраны решения с определенной спиральностью (проекцией спина на направление импульса) $\lambda= \pm 1 / 2$. В представлении Дирака-Паули (в к-ром $\gamma^{0}$ диагональна) эти решения имеют вид:



\[

\left.\begin{array}{c}

u_{\lambda}(p)=N\left(\begin{array}{l}

v_{\lambda} \\

\frac{2 \lambda|\boldsymbol{p}|}{p_{0}+m} v_{\lambda}

\end{array}\right) ; \quad \mathscr{E}=p_{0}, \quad \lambda= \pm 1 / 2, \\

u_{\lambda}(-p)=N\left(\begin{array}{l}

\frac{-2 \lambda|\boldsymbol{p}|}{p_{0}+m} v_{\lambda} \\

v_{\lambda}

\end{array}\right) ; \quad \mathscr{E}=-p_{0}, \quad \lambda= \pm 1 / 2 .

\end{array}\right\}

\]



Здесь $v_{\lambda}-2$-компонентный спинор, удовлетворяющий уравнению



\[

\frac{1}{2} \sigma \boldsymbol{n} v_{\lambda}=\lambda v_{\lambda},

\]



где $\boldsymbol{n}=\boldsymbol{p} /|\boldsymbol{p}|, \sigma\left(\sigma_{1}, \sigma_{2}, \sigma_{3}\right)-$ Паули матрицы, а множитель $N$ определяется нормировкой спинора $u( \pm p)$. Используются следуюшие нормировки (для каждого значения $\lambda$ ):



а) $(u( \pm p))^{+} u( \pm p)=1, \quad N=\sqrt{\left(p_{0}+m\right) / 2 p_{0}}$



б) $\bar{u}( \pm p) u( \pm p)= \pm 1, \quad N=\sqrt{\left(p_{0}+m\right) / 2 m}$,



в) $\bar{u}( \pm p) \gamma^{0} u( \pm p)=2 p_{0}, \quad N=\sqrt{p_{0}+m}$,



при этом $v^{+} v=1$.



Для $m=0$ решения свободного Д.у. являются собственными функциями матрицы $\gamma^{5}=-i \gamma^{0} \gamma^{1} \gamma^{2} \gamma^{3}$ :



\[

\gamma^{5} u_{\lambda}( \pm p)=\mp 2 \lambda u_{\lambda}( \pm p) .

\]



В матричные элементы процессов со слабым взаимодействием спиноры $u_{\lambda}(p)$, описывающие нейтрино, входят в виде $\frac{1}{2}\left(1+\gamma^{5}\right) u_{\lambda}(p)$. Если масса нейтрино равна нулю, то



\[

\begin{gathered}

\frac{1}{2}\left(1+\gamma^{5}\right) u_{\lambda}(p)=0 \quad \text { при } \lambda=1 / 2, \\

\frac{1}{2}\left(1+\gamma^{5}\right) u_{\lambda}(p)=u_{\lambda}(p) \quad \text { при } \lambda=-1 / 2,

\end{gathered}

\]



то есть спиральность нейтрино равна $-1 / 2$. Частице с отрицательной энергией соответствует антинейтрино (см. ниже), его спиральность равна $+1 / 2$.



В нерелятивистском случае $\beta=|\boldsymbol{p}| / p_{0} \ll 1$ (в системе СГС $\beta=v / c$, где $v-$ скорость частицы), и спиноры $u_{\lambda}( \pm p)$ с точностью до линейных по $\beta$ членов даются выражениями:



\[

u_{\lambda}(p)=N\left(\begin{array}{c}

v_{\lambda} \\

\beta \lambda v_{\lambda}

\end{array}\right), \quad u_{\lambda}(-p)=N\left(\begin{array}{c}

-\beta \lambda v_{\lambda} \\

v_{\lambda}

\end{array}\right) .

\]



Отсюда следует, что для нерелятивистской частицы «нижние» («верхние») компоненты решений Д. у. с положительной (отрицательной) энергией много меньше «верхних» («нижних») компонент.





\end{document}